\documentclass[11pt]{article}

\usepackage[a4paper,margin=1in]{geometry}
\usepackage[T1]{fontenc}
\usepackage{lmodern}
\usepackage{microtype}
\usepackage{amsmath,amssymb}
\usepackage{xcolor}
\usepackage{hyperref}
\usepackage{enumitem}
\usepackage{xurl}

\hypersetup{
  colorlinks=true,
  linkcolor=blue!50!black,
  urlcolor=blue!50!black,
  citecolor=blue!50!black,
  pdftitle={Formalizing Kraft's Inequality in Lean 4},
  pdfauthor={Elazar Gershuni}
}

\title{Formalizing Kraft's Inequality in Lean 4\\
  \large Final project summary}
\author{Elazar Gershuni\\
  \small Course 2360124, \emph{Proving theorems in Lean}\\[-2pt]
  \small Winter 2025/2026, Prof.\ Yuval Filmus, Technion}
\date{April 19, 2026}

% Concrete references (commit hashes, file paths, line numbers) are
% intentionally kept out of the rendered prose to stay readable. They live in
% the project's git log and in the companion notes next to this file.

\begin{document}
\maketitle

\begin{abstract}
\noindent
This report summarises a formalization project carried out for the course
\emph{Proving theorems in Lean}. The assigned task was to formalize the proof of Kraft's inequality; the scope ended up
expanding to arbitrary finite alphabets, an abstract monoid-level
formulation with nonnegative-real weights, an infinite-code variant, the
source coding lower bound via the Gibbs inequality, and a Shannon--Fano
application. Two small contributions were merged upstream, one to
\textsf{Mathlib} and one to \textsf{Lean~4} itself. This report
summarizes what was formalized, how the proofs are restructured relative
to the original, which pieces look simple in math but turn out tricky in
Lean, and the (uneven) experience of working with AI assistants.
\end{abstract}

\paragraph{Links.}
\begin{itemize}[leftmargin=1.5em,itemsep=0pt]
  \item Repository and \texttt{README.md}: \url{https://github.com/elazarg/kraft}
  \item Mathematical specification (\texttt{kraft.tex}):
    \url{https://github.com/elazarg/kraft/blob/main/kraft.tex}
  \item \textsf{Mathlib} PR~\#34108 (merged):
    \url{https://github.com/leanprover-community/mathlib4/pull/34108}
  \item \textsf{Lean~4} PR~\#12108 (merged):
    \url{https://github.com/leanprover/lean4/pull/12108}
\end{itemize}

\medskip\noindent
Project duration: 2025-12-07 to 2026-01-23 (about seven weeks,
$\approx$150 commits). Toolchain: \textsf{Lean~4.26.0} with
\textsf{Mathlib~v4.26.0}. License: Apache~2.0.

\section{The task and the original proofs}

The task was to formalize the contents of the course note \texttt{kraft.tex}
in Lean~4 with \textsf{Mathlib}. The note fixes the binary alphabet
$\{0,1\}$ and states four results:

\begin{itemize}[leftmargin=1.5em,itemsep=1pt]
  \item \textbf{Kraft's inequality.} For a finite prefix-free $S \subseteq
    \{0,1\}^*$, $\ \sum_{w \in S} 2^{-|w|} \le 1$.
  \item \textbf{Converse.} For any $\ell\colon I \to \mathbb{N}$ on a finite
    $I$ with $\sum_i 2^{-\ell(i)} \le 1$, there is an injective
    $w\colon I \to \{0,1\}^*$ whose image is prefix-free and
    $|w(i)| = \ell(i)$.
  \item \textbf{Kraft--McMillan.} The same bound $\sum 2^{-|w|} \le 1$
    holds under the weaker hypothesis that $S$ is \emph{uniquely decodable}
    -- i.e.\ concatenation is injective on lists of codewords.
  \item \textbf{Prefix-free $\Rightarrow$ uniquely decodable}, provided
    $\varepsilon \notin S$.
\end{itemize}

\noindent
The mathematical proofs, in broad strokes -- just enough to follow the rest
of this report; see \texttt{kraft.tex} for details.

\paragraph{Kraft (counting).} Let $\ell := \max_{w \in S} |w|$. For each
$w \in S$ the extension set $E(w) = \{ wz : |z| = \ell - |w| \}$ sits inside
$\{0,1\}^\ell$ and has size $2^{\ell - |w|}$. Prefix-freeness makes the
$E(w)$ pairwise disjoint, so $2^\ell \ge \sum_{w \in S} 2^{\ell - |w|}$;
divide by $2^\ell$.

\paragraph{Converse.} Induct on $\max \ell(i)$. The heart is an auxiliary
lemma: if the $\ell(i)$ are sorted nondecreasingly and
$\sum 2^{-\ell(i)} \ge 1$, some prefix of the sequence has sum
\emph{exactly}~$1$. A corollary gives a subset of weight
$\tfrac{1}{2}$ when the total weight is above $\tfrac{1}{2}$ and all
lengths are positive. The proof uses this to split $I$ into two parts
whose Kraft weights are both at most $\tfrac{1}{2}$; shift to
$\ell'(i) := \ell(i) - 1$; recurse on each side; prepend $0$ to the words
on the left and $1$ on the right.

\paragraph{Kraft--McMillan.} Let $C := \sum_{w \in S} 2^{-|w|}$. Then
\[
  C^r \;=\; \sum_{(w_1,\dots,w_r) \in S^r} 2^{-|w_1 \cdots w_r|}.
\]
Unique decodability makes $(w_1,\dots,w_r) \mapsto w_1 \cdots w_r$
injective, and the image lies in $\bigcup_{s=r}^{r\ell} \{0,1\}^s$, so
$C^r \le \sum_{s=r}^{r\ell} 2^s \cdot 2^{-s} = r\ell - r + 1$.
Take $r$-th roots and send $r \to \infty$: $C \le 1$.

\paragraph{Prefix-free $\Rightarrow$ uniquely decodable.} Induct on $|x|$.
If $x = w_1 \cdots w_r = z_1 \cdots z_s$ with $|w_1| \le |z_1|$, then $w_1$
is a prefix of $z_1$, hence $w_1 = z_1$ by prefix-freeness; strip and
recurse. (The base case forces the side condition $\varepsilon \notin S$:
otherwise $\varepsilon = \varepsilon\varepsilon$ is two decodings of the
same string.)

\section{The formalization}
\label{sec:formalization}

All four theorems are formalized, considerably generalised. None of the
following additions appears in the course note; each emerged naturally
while proving the originals.

\begin{itemize}[leftmargin=1.5em,itemsep=1pt]
  \item \textbf{Arbitrary finite alphabet of size $D \ge 2$.}
  \item \textbf{Infinite codes.} Kraft's inequality and the converse are
    proved for possibly-infinite index sets.
  \item \textbf{Abstract monoid formulation.} Kraft--McMillan is proved
    once for any graded monoid $M$ satisfying a $D$-ary growth bound, with
    a multiplicative weight
    $\mu\colon M \to \mathbb{R}_{\ge 0}$ dominated by
    $D^{-\mathrm{cost}(x)}$, using a monoid-level injectivity hypothesis
    in place of unique decodability.
  \item \textbf{Real-valued weights.} The abstract theorem is stated over
    $\mathbb{R}_{\ge 0}$, with counting and limits separated.
  \item \textbf{Source coding lower bound.} For any injective uniquely
    decodable code $w$ and strictly positive distribution $p$,
    \[
      H_D(p) := -\!\sum_i p_i \log_D p_i \;\le\; \mathbb{E}_{i \sim p}\bigl[|w(i)|\bigr].
    \]
  \item \textbf{Shannon--Fano.} For any strictly positive finite
    distribution there is a prefix-free code with
    $\mathbb{E}[|w|] < H_D(p) + 1$, obtained from the lengths
    $\ell(i) = \lceil -\log_D p_i \rceil$.
  \item \textbf{Two upstream PRs} (\S\ref{sec:upstream}).
\end{itemize}

\medskip
\noindent
The formalization reuses many of the mathematical ideas of the original
proofs -- especially the counting behind Kraft and the $r$-fold power trick
for McMillan -- but partitions the work so that each reusable idea appears
once. The converse is the main exception: the final Lean proof uses a
canonical interval construction rather than the note's recursive halving
argument.

\paragraph{Kraft mostly becomes a corollary of Kraft--McMillan.} The note
proves Kraft's inequality first, as a standalone theorem, before
introducing McMillan. In the formalization, the standalone counting proof
is absent. Instead:
\begin{enumerate}[leftmargin=1.5em,itemsep=0pt]
  \item ``Prefix-free $\Rightarrow$ uniquely decodable'' is proved as its
    own theorem under the necessary side condition $\varepsilon \notin S$.
  \item Kraft--McMillan is proved for uniquely decodable codes.
  \item Kraft's inequality for prefix-free codes is then the trivial
    $\{\varepsilon\}$ case plus the Kraft--McMillan corollary.
\end{enumerate}
The counting argument of the original Kraft proof is still there -- but it
appears once, as part of McMillan, not twice.

\paragraph{McMillan is split along $\mathbb{N}$ vs.\ $\mathbb{R}$.} The
original does counting and the $r \to \infty$ limit in a single paragraph.
The formalization separates them into two statements at different levels
of abstraction:

\begin{enumerate}[leftmargin=1.5em,itemsep=0pt]
  \item \textbf{Counting (pure $\mathbb{N}$, monoid-level).} For a finite
    $S$ in a graded monoid satisfying a $D$-ary growth bound, if $r$-fold
    products are injective, the scaled tuple-counting sum is bounded by
    \[
      \sum_{w \in S^r}
        D^{N-\mathrm{cost}(w_1\cdots w_r)}
      \;\le\; (N+1)D^N,
      \qquad N := r \cdot \max_{x \in S} \mathrm{cost}(x).
    \]
    This is the original's counting argument with the alphabet, codes, and
    real numbers all stripped away.
  \item \textbf{Limit (pure $\mathbb{R}_{\ge 0}$).} If $K^r$ is bounded by
    a linear function of $r$ for every $r$, then $K \le 1$. This is where
    the formalization introduces a multiplicative weight
    $\mu\colon M \to \mathbb{R}_{\ge 0}$ dominated by
    $D^{-\mathrm{cost}(x)}$.
  \item \textbf{Specialisation to lists.} About ten lines: list length is
    additive, there are at most $D^s$ strings of length $s$, unique
    decodability means injectivity of $r$-fold concatenation.
\end{enumerate}

\paragraph{The same machine drives the source coding bound.} The source
coding lower bound uses the same abstract machinery with a different weight
on lists: the normalised Kraft weight $q_i := D^{-|w(i)|} / K$. Applying
the Gibbs inequality $0 \le \sum p_i \log(p_i / q_i)$ to $p$ and $q$ --
itself a statement about two probability distributions, not about codes --
gives, after cancelling $\log K \le 0$,
\[
  H_D(p) \;\le\; \mathbb{E}_p \bigl[|w|\bigr].
\]
Shannon--Fano then falls out of the \emph{converse} by choosing
$\ell(i) = \lceil -\log_D p_i \rceil$: a small calculation shows
$\sum D^{-\ell(i)} \le 1$, the converse supplies a code, and
$\mathbb{E}[|w|] < H_D(p) + 1$ because $\lceil x \rceil < x + 1$.

\paragraph{The converse stays separate, because it isn't a counting argument.}
It is algorithmic: the proof \emph{produces} a code. Here the
formalization deliberately diverges from the recursive halving proof in
the note. It sorts the indices by nondecreasing length, assigns each index
a $D$-adic interval start
\[
  A_n/D^{\ell_n}
  = \sum_{k<n} D^{-\ell_k},
\]
and encodes $A_n$ as a fixed-width base-$D$ word of length $\ell_n$.
The strict prefix-sum bound keeps $A_n < D^{\ell_n}$, while a division
separation lemma proves prefix-freeness. Finite, infinite, and arbitrary
alphabet cases are then wrappers around this construction.

\section{Math vs Lean}
\label{sec:mathvslean}

A tour of steps where the original proof is a single sentence and the
formalization has to set up machinery; plus a few \textsf{Mathlib}
concepts a reader of the course note might not recognise.

\paragraph{``Arrange the lengths in order.''} The converse reorders
silently on paper. For the subset-sum lemma, and also for the formalized
interval construction, the useful order is nondecreasing length -- shorter
codewords first. Lean wants the ordering as an explicit object: a monotone
bijection
$\mathrm{Fin}\,n \simeq I$ for finite $I$, or $\mathbb{N} \simeq I$ for
infinite $I$ once the sequence is known summable. Two helper equivalences
isolate the existence of the ordering from its use, so the main proof can
quote them without repeating the construction.

\paragraph{``Take $r$-th roots and let $r \to \infty$.''} One sentence in
the original; avoided in Lean, because real-valued $r$-th roots in
$\mathbb{R}_{\ge 0}$ are awkward to manipulate. The contradiction
argument sketched in \S\ref{sec:formalization} replaces them: the
relevant \textsf{Mathlib} fact is $r \cdot c^r \to 0$ whenever $|c| < 1$,
and setting $c = 1/K$ closes the gap without ever computing a root.

\paragraph{Infinite codes, and what \texttt{tsum} is.} \textsf{Mathlib}
has two summation notations. For a finite $S$, \texttt{$\sum$ i $\in$ S, f i}
is the ordinary finite sum. For a possibly-infinite index type,
\texttt{$\sum$' i, f i} (read ``\texttt{tsum}'', for ``topological sum'')
denotes the limit of finite partial sums \emph{when the family is
summable} -- the predicate \texttt{Summable} -- and is defined to be $0$
otherwise, for convenience. So ``Kraft's inequality for infinite
prefix-free codes'' really says two things: the family $i \mapsto
D^{-|w_i|}$ is summable, and its \texttt{tsum} is at most $1$. The
formalization proves both together by bounding every finite partial sum
by $1$ via the finite Kraft theorem -- a pattern that applies to any
``convergent by finite subsums'' statement.

\paragraph{Truncated subtraction in $\mathbb{N}$.} The original writes
$N - |w|$ freely, assuming $|w| \le N$ from context. Lean's
$\mathbb{N}$-subtraction is truncated to $0$ on the overflow case, so any
algebraic manipulation of an expression like $D^{N-|w|}$ carries a side
condition $|w| \le N$. A midway refactor rewrote the counting bound so
that invariants use only additive identities, confining subtraction to
one boundary lemma.

\paragraph{$\mathbb{R}_{\ge 0}$ is its own type.} \textsf{Mathlib}'s
\texttt{NNReal} (denoted $\mathbb{R}_{\ge 0}$) is implemented as the
subtype of real numbers satisfying $0 \le x$, but it is still a separate
Lean type with its own arithmetic instances. Positivity is baked in, so
$0 \le x$ hypotheses do not float around, and $(\cdot)^{-1}$ behaves
sensibly at $0$. The cost is that casts between $\mathbb{R}_{\ge 0}$ and
$\mathbb{R}$ must be done deliberately. The abstract Kraft theorem is
stated in $\mathbb{R}_{\ge 0}$; a small wrapper provides the
$\mathbb{R}$-valued version that the list and source-coding theorems
consume.

\paragraph{KL divergence and the $0 \log 0$ convention.} The textbook
source coding lower bound uses the Gibbs inequality
\[
  \sum_i p_i \log(p_i/q_i) \ge 0,
\]
which in prose relies on $0 \log 0 = 0$. \textsf{Mathlib} bakes the
convention into \texttt{Real.negMulLog}; the current theorem assumes
strictly positive probabilities, but this is the convention needed for the
usual zero-probability extension. The formalization goes further: it
re-expresses $p$ and $q$ as finite measures, invokes \textsf{Mathlib}'s
KL-divergence API \texttt{klDiv}, and transports its non-negativity
theorem back to the finite-sum form.
The price is that the KL-based route proves and uses the absolute
continuity $p \ll q$, and the finite-sum/log transport is currently
written for strictly positive $p_i$ and $q_i$. The normalized Kraft weights
give $q_i>0$ automatically; the public theorem therefore assumes
$p_i > 0$. A relaxation via support truncation is noted in the code as
future work.

\section{Upstream contributions}
\label{sec:upstream}

\paragraph{\textsf{Mathlib} PR~\#34108 (merged, $\approx$\,240 lines).}
\emph{feat(InformationTheory/Coding): Kraft--McMillan inequality for
uniquely decodable codes.} Establishes new
\path{Mathlib.InformationTheory.Coding} modules in the
\texttt{InformationTheory} namespace, adding the
definition of uniquely decodable codes, the impossibility of containing
the empty word, and the Kraft--McMillan inequality itself. A small
supporting lemma about filtered sums was added to \textsf{Mathlib}'s
\texttt{BigOperators} infrastructure. Much of the work on the PR was golfing, which actually turned the proof into a cleaner one.

\paragraph{\textsf{Lean~4} PR~\#12108 (merged, $8$ lines).} Adds two
small utility lemmas stating that, for injective $f$, $l_1.\mathrm{map}\,f$
is a prefix (respectively a suffix) of $l_2.\mathrm{map}\,f$ if and only
if $l_1$ is a prefix (suffix) of $l_2$. This is the exact fact needed
when lifting a code over one alphabet to one over a larger alphabet via
an embedding -- small, but reusable enough to upstream.

\section{AI interaction}
\label{sec:ai}

Three assistants contributed, with very uneven weights. \textbf{ChatGPT}
(versions 5.2 and 5.3) was the primary tool, driving most day-to-day proof
drafting and refactoring. \textbf{Gemini~3} was a secondary opinion,
consulted when ChatGPT stalled; it sometimes agreed with ChatGPT in a way
that turned out to be wrong. \textbf{Aristotle} (Harmonic) was used once,
for one critical step, with outsized impact. A coding agent was essentially
not used during the project, except for a little polish near the end.

\subsection*{The Aristotle episode}

One mid-January morning I was stuck on the auxiliary subset-sum lemma then
underpinning the converse: given sorted lengths with $\sum 2^{-\ell_i} \ge
1$, find a prefix whose sum is \emph{exactly}~$1$. I asked Aristotle for a
proof and, in parallel, ChatGPT and Gemini for their opinions on it. Both
ChatGPT and Gemini insisted Aristotle's approach was the wrong direction
and should be discarded in favor of their own preferred route. I
committed Aristotle's proof anyway, cleaned it up, and discovered that
the ``wrong'' proof needed only small gaps filled. By the end of the same
day, that version of the converse was closed. The final code later moved
to the interval construction described above, so this episode survives as
project history rather than as a lemma in the current files.

The lesson: \emph{``both mainstream models agree this is wrong'' is not
authoritative}. Two frontier LLMs dismissing a proof is at most weak
evidence. If the goals close, the proof is right.

\subsection*{Where AI was strong}

\begin{itemize}[leftmargin=1.5em,itemsep=0pt]
  \item Bookkeeping-heavy algebra. Rearranging
    $\bigl(\sum_x \mu(x)\bigr)^r$ into $\sum_{w \in S^r} \mu(w_1 \cdots w_r)$
    and chasing cancellations like $D^{N-c}(D^N)^{-1} = D^{-c}$ are
    tedious, mechanical, and suit an LLM well. Drafts were usually right
    after a revision or two.
  \item Golfing. ``Can this be shorter?'' often produced a real
    simplification, especially after a first working version.
  \item Naming. The abstractions named \texttt{WeightModel} and
    \texttt{ExpBounded} carry the generalized proof; both were
    AI-suggested and survived external review.
  \item Boilerplate. Case splits and type-transport plumbing -- the
    finite/infinite split in the converse, lifting a code over
    \texttt{Fin}~$D$ into one over a general alphabet -- were easy wins.
\end{itemize}

\subsection*{Where AI was weak, and the cleanup it triggered}

\begin{itemize}[leftmargin=1.5em,itemsep=0pt]
  \item \textbf{Hallucinated lemma names.} Pervasive. Any named lemma that
    looked plausible had to be grep-verified before it was trusted. The
    discipline that emerged: always search (\texttt{loogle},
    \texttt{leansearch}, or plain grep in \textsf{Mathlib}) before using
    a name.
    
  \item \textbf{Stale conventions.} Old lemma names, old syntax (Lean 3). Old conventions (``avoid golfing'').

  \item \textbf{Makes it easy to lose track.} It's tempting to let the AI write the code. And it works - until it doesn't, and then you're at a loss.
\end{itemize}

\subsection*{Division of labour and external review}

Architecture decisions (when to generalize, file layout, what to upstream)
and every final-commit decision were mine. The AI drafted short tactic
blocks and proposed refactors; each block was read in full before being
committed, and many commits titled ``cleanup'' or ``simplify'' are
follow-ups on AI-produced proofs.

The external check came from the \textsf{Mathlib} PR (\S\ref{sec:upstream}).
Across three rounds of review, the only requests were stylistic: calc-block
indentation, inlining of intermediate types, and a reformulation to an
indexed-family form. No mathematical revisions were requested.

\section{Advice for future students}

\begin{enumerate}[leftmargin=1.5em,itemsep=1pt]
  \item \textbf{Budget a generalisation pass.} The first end-to-end proof
    is rarely the one you keep. Plan to revisit the main theorem once in
    a more abstract setting; it usually simplifies the concrete case too.
  \item \textbf{Keep $\mathbb{N}$ and $\mathbb{R}$ apart.} Subtraction
    and division break differently. Splitting the natural-number counting
    bound from the real-valued limit argument was a late but decisive
    refactor.
  \item \textbf{Use automation where agents are stuck.} Try
    \texttt{grind}, \texttt{simp\_all}, \texttt{aesop?}, and
    \texttt{exact?}; they often close goals that an agent keeps circling.
  \item \textbf{Search before guessing lemma names.} One edit-compile
    cycle saved per correct guess. A language-server MCP like
    \texttt{lean-lsp-mcp} lets the assistant query goal state, look up
    signatures, and search local declarations instead of inventing names.
  \item \textbf{Don't be surprised if an important utility does not exist.} 
  \item \textbf{Let \textsf{Mathlib} do the heavy lifting.} If a concept
    has a classical name, grep for it before rolling your own. The
    source coding lower bound, for instance, fits into \textsf{Mathlib}'s
    existing KL-divergence API almost for free (by sheer luck - there's not much information theory in \textsf{Mathlib}).
  \item \textbf{In a small project, experiment freely; allow complete rewrites from an LLM agent.} Sunk costs are not large, and a clean-slate pass can resolve an architectural knot faster than a sequence of local edits. It also  tends to expose the real structural lemmas -- the ones the next draft will factor out.
  \item \textbf{Listen to the proofs.} they may contain interesting abstractions hidden within.
  \item \textbf{Never trust a definition written by an agent.} An agent
    asked to \emph{prove} a claim will quietly \emph{weaken} it when the
    proof fails: drop a hypothesis, strengthen a premise. Definitions are worse: a slightly shifted definition can
    make the statement trivially true in the wrong way. Read every
    definition the agent introduces, by hand, against your own mental
    statement of what it should mean.
\end{enumerate}

\end{document}
